\documentclass[12pt]{report}
\usepackage[top=2.3cm,bottom=2.5cm,left=2.5cm,right=2.5cm]{geometry}
\usepackage[T1]{fontenc}
\usepackage{mathpazo}
\linespread{1.03}
\usepackage{microtype}
\usepackage{graphicx}
\usepackage{booktabs}
\usepackage{enumitem}
\setlist{itemsep=0.15em,topsep=0.35em,parsep=0pt}
\usepackage{array}
\usepackage{float}
\usepackage{longtable}
\usepackage{tabularx}
\usepackage{amsmath}
\usepackage{amssymb}
\setlength{\parindent}{1.4em}
\setlength{\parskip}{0pt}
\usepackage[hidelinks]{hyperref}

\begin{document}
\pagestyle{plain}
\setcounter{page}{4}
\setcounter{chapter}{2}
\chapter{Memory Management}

\section{Memory Management}

Memory management is the core responsibility of an operating system that involves allocating, tracking, and reclaiming main memory for processes. This chapter introduces the concepts of virtual memory and the mechanisms that support it.

The abstraction of virtual memory allows each process to view the available memory as a contiguous address space, independent of the physical RAM installed in the machine. The operating system maintains data structures that translate virtual addresses used by programs into physical addresses in RAM or, when necessary, into secondary storage.

\subsection{Virtual Memory Concepts}
Virtual memory provides isolation and protection between processes while extending the address space beyond the limits of physical memory. It enables features such as memory‑mapped files and on‑demand loading of code and data.

\subsection{Paging Mechanism}
Paging divides both virtual and physical memory into fixed‑size blocks called pages and frames respectively. The operating system uses a page table to keep track of the mapping between a process’s pages and the available frames.

\begin{verbatim}
+-------------------+      +-------------------+
|   Process A       |      |   Physical Memory |
| Code | Data |Heap |      | Frame 0 | Frame 1 |
+-------------------+      +-------------------+
|   Process B       |      | Frame 2 | Frame 3 |
+-------------------+      +-------------------+
|   Free Frames     |      |      ...          |
+-------------------+
\end{verbatim}

\subsection{Page Fault Handling}
When a process accesses a page that is not resident in physical memory, a page fault occurs. The operating system handles the fault by allocating a frame, reading the required page from disk, updating the page table, and restarting the interrupted instruction.

\begin{verbatim}
   +-------------------+    Page Fault    +-------------------+
   |  CPU accesses     |  -------------->|  Page Table Lookup|
   |  memory address   |                 |  (not present)    |
   +-------------------+                 +-------------------+
                                                |
                                                v
                                        +-------------------+
                                        |  Page Fault       |
                                        |  Exception        |
                                        +-------------------+
                                                |
                                                v
                                        +-------------------+
                                        |  Page Fault       |
                                        |  Handler          |
                                        +-------------------+
                                1. Find free frame / evict
                                2. Load page from disk
                                3. Update page table
                                4. Restart faulting instruction
\end{verbatim}

\subsection{Comparison of Physical and Virtual Memory}
\begin{table}[htbp]
\centering
\begin{tabular}{ll}
\toprule
Aspect & Description \\
\midrule
Physical Memory & The actual RAM installed in the computer; limited in size and directly accessible by the CPU. \\
Virtual Memory & An abstraction that provides each process with an address space larger than physical memory, using disk as an extension. \\
\bottomrule
\end{tabular}
\caption{Comparison of Physical and Virtual Memory}
\end{table}

\subsection{Swap Space}
Swap space is a designated area on secondary storage that the operating system uses to store pages that are not currently in physical memory. It allows the system to overcommit memory and to free frames for other processes.

\subsection{Summary}
Memory management combines allocation strategies, protection mechanisms, and swapping policies to present a coherent virtual address space to processes. The next chapter will explore how files are organized and accessed through the file system and I/O subsystem.

\section*{Tutorial Questions}
\begin{enumerate}
\item Explain the difference between physical memory and virtual memory. Provide one advantage of using virtual memory for process isolation.
\item Describe how a page table maps virtual pages to physical frames. Include a brief example with two processes.
\item What is a page fault and what steps does the operating system take to resolve it?
\item Compare the concepts of a frame and a page. Why must they be of the same size?
\item Discuss the purpose of swap space. How does it affect system performance?
\item A process attempts to access a page that is not in memory, causing a page fault. Explain the sequence of events that follow, from detection to resumption of the process.
\item Identify two scenarios where demand paging improves overall system throughput.
\item Outline the conditions under which a page replacement algorithm must be invoked.
\end{enumerate}

\end{document}
